% TOP_PROBLEM: {"id": 42, "title": "Whitehead Problem", "category_id": 10, "category": {"id": 10, "name": "set_theory", "display_name": "Set Theory", "description": "Foundations of mathematics, infinite sets, and cardinality.", "slug": "set-theory", "order_index": 10, "created_at": "2026-07-31T15:26:25.671Z"}, "status": "open"}
% FIRST_POSTED: null
% ENTRY_KIND: statement_only
% Imported from: batch01.tex; UnsolvedMath ID 42
\documentclass[11pt]{article}
\usepackage[margin=25mm]{geometry}
\usepackage{fontspec}
\setmainfont{Latin Modern Roman}
\usepackage{amsmath,amssymb,mathtools}
\usepackage{unicode-math}
\setmathfont{Latin Modern Math}
\usepackage{xurl,hyperref}
\hypersetup{colorlinks=true,urlcolor=blue}
\setcounter{secnumdepth}{0}
\setlength{\parindent}{0pt}
\setlength{\parskip}{5pt}
\setlength{\emergencystretch}{3em}
\newcommand{\sref}[2]{\hyperlink{source:#1:#2}{[S#2]}}
\newcommand{\cataloguescope}[1]{\paragraph{Recorded target.}#1}
\begin{document}
% UnsolvedMath ID 42; source catalogue code SET-003
\setcounter{section}{41}
\section{Whitehead Problem}\label{42}
{\small\sffamily UnsolvedMath ID 42 · Set Theory}
\subsection{Definitions and mathematical statement}

\paragraph{Shared notation.}$\mathbb N=\{1,2,\ldots\}$, and $\mathbb Z,\mathbb Q,\mathbb R,\mathbb C$ denote the integers, rationals, reals and complex numbers. Any other conventions are specified locally.


All groups here are ordinary abelian groups, with no topology and no restriction on cardinality. The group $\operatorname{Ext}^1_{\mathbb Z}(A,\mathbb Z)$ classifies extensions of $A$ by $\mathbb Z$ up to equivalence. Its vanishing means that every short exact sequence of abelian groups
\[
 0\longrightarrow\mathbb Z\xrightarrow{i}B\xrightarrow{\pi}A\longrightarrow0
\]
splits: there is a homomorphism $s:A\to B$ with $\pi\circ s=\operatorname{id}_A$. Short exactness means $i$ is injective, $\pi$ is surjective, and $\operatorname{im}(i)=\ker(\pi)$. A group is free abelian if it is isomorphic to $\bigoplus_{j\in I}\mathbb Z$ for some set $I$; elements of this direct sum have finite support.
\paragraph{Statement recorded in the supplied source.}
Is every abelian group $A$ satisfying
\[
 \operatorname{Ext}^1_{\mathbb Z}(A,\mathbb Z)=0
\]
free abelian? Equivalently, does splitting of every extension by $\mathbb Z$ force $A\cong\bigoplus_{j\in I}\mathbb Z$ for some set $I$? \sref{42}{2}
This historical question is independent of ZFC, under the usual consistency assumption. Shelah's positive result under the axiom of constructibility and a consistent negative result are recalled in Source~S2, Theorem 1.1. The assertion is therefore not a universal theorem or refutation of ZFC. \sref{42}{2}
\subsection{Short English statement}
If every abelian-group extension of a group by the integers splits, must that group have a basis over the integers? The unrestricted question is known to be independent of ZFC.
\subsection{Sources}
\begin{description}
\item[\hypertarget{source:42:1}{S1}] ULAM.AI and the UnsolvedMath Contributors, UnsolvedMath: A Curated Collection of Open Mathematics Problems, record 42 (SET-003). CC BY 4.0. \url{https://huggingface.co/datasets/ulamai/UnsolvedMath}.
\item[\hypertarget{source:42:2}{S2}] Jeffrey Bergfalk, Chris Lambie-Hanson and Jan Šaroch, Whitehead’s Problem and Condensed Mathematics, arXiv:2312.09122v5 (2025), Section 1, original ordinary-group question and Theorem 1.1. \url{https://arxiv.org/pdf/2312.09122}.
\end{description}
\label{end:42}
\end{document}
